\largerpage
\scp{Internal structure}{06-02}
Within the \tsc{Aru} subgroup there is a \tsc{Nuclear Aru} node
which further contains a \tsc{South Aru} node.
The \tsc{Nuclear Aru} node is defined by *s > \it{h} > {\0},
which occurs in all languages except Ujir.
This leaves Ujir as a family level isolate within \tsc{Aru}.
West Kola has *s > \it{h} while other languages have *s > {\0},
probably through intermediate **h.
Although some languages appear to have *s > \it{y},
this is actually due to glide insertion after loss of *s,
or a shift of final *i > \it{y} after loss of intervocalic *s.
Examples of PMP *s are given in \trf{06-13}.

\begin{table}\tcp{\tsc{Aru} *s}{06-13}
\begin{adjustbox}{width=\textwidth}
	\begin{threeparttable}\st{2pt}
	\begin{tabular}{lllllllll}\lsptoprule
local gl.	&	{\hr}snake	&	{\hr}they	&	{\hr}one 	&	{\hr}what	&	{\hr}sago tree	&	{\hr}sea	&	{\hr}bow	&	{\hr}paddle\\
PMP	&	\hr*\tbr{s}awa\su{†}	&	\hr*\tbr{s}i-ida	&	{\hr}?\hr*ə\tbr{s}a	&		&		&	\hr*ta\tbr{s}ik	&	\hr*bu\tbr{s}uR	&	\hr*bəR\tbr{s}ay\\
Proto-Aru	&	**\tbr{s}awa	&	**\tbr{s}ira	&	**\tbr{s}etu	&	**\tbr{s}a\su{‡}	&	**\tbr{s}akaw	&	**ta\tbr{s}i	&	**pu\tbr{s}uR	&	**peR\tbr{s}ay\\ \midrule
Ujir	&	{\hr\hr}\tbr{s}awa	&	{\hr\hr}\tbr{s}ira	&	{\hr\hr}\tbr{s}et	&	{\hr\hr}\tbr{s}a	&	{\hr\hr}\tbr{s}akaw	&	{\hr\hr}te\tbr{s}	&	{\hr\hr}ɸu\tbr{s}i	&	{\hr\hr}ɸe\tbr{s}a\\
Kompane	&	{\hr\hr}yaa	&		&	{\hr\hr}etu	&	{\hr\hr}ya	&	{\hr\hr}akaw	&		&		&	\\
West Kola	&	{\hr\hr}\tbr{h}awa	&		&	{\hr\hr}\tbr{h}ot	&	{\hr\hr}en \tbr{h}a	&	{\hr\hr}\tbr{h}akaw	&		&		&	\\
East Kola	&	{\hr\hr}awa	&	{\hr\hr}iri	&	{\hr\hr}ot	&	{\hr\hr}ye	&	{\hr\hr}akaw	&	{\hr\hr}tai rei	&	{\hr\hr}ɸuuh	&	{\hr\hr}ɸe\tbr{h}a\su{ǁ}\\
Lola	&	{\hr\hr}yawa	&		&	{\hr\hr}ye	&	{\hr\hr}menayá	&	{\hr\hr}yaʔu	&		&		&	\\
Lorang	&	{\hr\hr}yagwa	&		&	{\hr\hr}yetu	&	{\hr\hr}ya	&	{\hr\hr}yaka	&		&		&	\\
Dobel	&	{\hr\hr}yakwa	&	{\hr\hr}ʔiri, yiri	&	{\hr\hr}e‹t›tú	&	{\hr\hr}ya	&	{\hr\hr}yaʔu	&	{\hr\hr}tay	&	{\hr\hr}ɸiyar	&	{\hr\hr}ɸeri\\
Koba	&	{\hr\hr}yáʍa	&		&	{\hr\hr}ekatú	&	{\hr\hr}ya	&	{\hr\hr}yaka	&	{\hr\hr}tay	&		&	\\
Barakai	&	\hr\cg{(ʤel)}	&		&	{\hr\hr}etí	&	\hr\cg{(woy)}	&	{\hr\hr}aga	&		&		&	\\
Karey	&	\hr\cg{(ʤel)}	&		&	{\hr\hr}etú	&	\hr\cg{(oín)}	&	{\hr\hr}agaw	&		&		&	\\
Batuley	&	{\hr\hr}aw	&	{\hr\hr}id	&	{\hr\hr}et	&	\hr\cg{(ula)}	&	{\hr\hr}aga	&	{\hr\hr}tay	&	{\hr\hr}ɸo(w)or	&	{\hr\hr}ɸer\\
Mariri	&		&		&	{\hr\hr}eʤ	&	\hr\cg{(totóo)}	&	{\hr\hr}ágao	&	{\hr\hr}tay	&		&	\\
Manombai	&		&	{\hr\hr}ira	&	{\hr\hr}etu	&	\hr\cg{(akarey)}	&	{\hr\hr}eka	&		&		&	\\
NE. Tarang.	&	{\hr\hr}yɔ	&	{\hr\hr}ira	&	{\hr\hr}ôt	&	{\hr\hr}ʤa	&	{\hr\hr}aka	&		&		&	\\
SE. Tarang.	&	{\hr\hr}iɔ	&	{\hr\hr}dir	&	{\hr\hr}ɛt-na	&	{\hr\hr}ʤa	&	{\hr\hr}agaw	&		&	{\hr\hr}ɸɔyar	&	\\
NW. Tarang.	&	{\hr\hr}yɔ, iɔ	&	{\hr\hr}ir	&	{\hr\hr}ôt	&	{\hr\hr}ya	&	{\hr\hr}aka	&	{\hr\hr}tay\su{Δ}	&	{\hr\hr}ɸɛar	&	{\hr\hr}ɸɛra\\
SW. Tarang.  	&	\hr\cg{(ʤɛla)}	&	{\hr\hr}êra	&	{\hr\hr}ôt	&	\hr\cg{(lebá)}	&	{\hr\hr}okaw	&	{\hr\hr}tay	&	{\hr\hr}ɸɛr	&	{\hr\hr}ɸera\\
		\lspbottomrule
	\end{tabular}
		\begin{tablenotes}\tnsize
			\item[†]	{PMP *sawa is reconstructed with the meaning ‘python’.}
			\item[‡]	{Likely cognates of **sa ‘what’ occur in \tsc{Timor-Babar} languages:
								Proto-Rote-Meto **saa, Helong \it{saa}, Tetun \it{saa} ‘what’ \citep[319]{Edwards2021},
								though these could be reflexes of PMP *sapa.}
			\item[Δ]	{Northwest Tarangan \it{tay} is from Juring village
								and \it{ɸɛra} ‘paddle’ from Rebi village.}
			\item[{ǁ}]	{Medial \it{h} in East Kola \it{ɸeha} ‘paddle’ is probably
								due to regular **R > \it{h},
								not *s > \it{h}.}
		\end{tablenotes}
	\end{threeparttable}
\end{adjustbox}
\end{table}

Most \tsc{Nuclear Aru} languages also share instances of *t >
**s, normally before front vowels.
However, this change is not complete in most languages,
and exceptions are common.
It probably spread between \tsc{Nuclear Aru} languages by diffusion,
or it may have begun at Proto-Nuclear Aru in a drift-like fashion,
with languages carrying it out to different extents.
Examples of *t > \it{t},
**s are given in \trf{06-14}.
%with instances of *t > **s shaded.
Some languages have subsequent **s > {\0}.
Glides are often subsequently inserted before vowel-initial words
that result from initial **s > {\0}.

\begin{table}\tcp{\tsc{Nuclear Aru} *t > \it{t}, \it{s}}{06-14}
\begin{adjustbox}{width=\textwidth}
	\begin{threeparttable}\st{2pt}
	\begin{tabular}{llllllll}\lsptoprule
local gl.	&	nine	&	wound	&	thick	&	three	&	plant (v.)\su{‡}	&	{\hr}we (incl.)	&	see\\
PMP	&		&		&	bamboo	&		&		&	\hr*kita	&	\\
Proto-Aru	&	**\tbb{t}era\su{†}	&	**\tbb{t}umbu	&	**\tbb{t}emal	&	**la\tbb{t}i	&	**-sa\tbb{t}i	&	**\tbb{t}i\tbb{t}a	&	**-\tbb{t}ewa\\ \midrule
Ujir	&	{\hr\hr}\tbb{t}era	&	{\hr\hr}\tbb{t}ubu	&	{\hr\hr}\tbb{t}emal	&	{\hr\hr}la\tbb{t}i	&	{\hr\hr}sa\tbb{t}	&	{\hr\hr}\tbb{t}ita	&	{\hr\hr}-\tbb{t}o\\
Kompane	&	{\hr\hr}\tbb{t}eri	&	{\hr\hr}\tbr{s}ubu	&		&	{\hr\hr}la\tbr{s}i	&		&		&	{\hr\hr}-\tbr{s}oo\\
East Kola	&	{\hr\hr}\tbb{t}era	&	{\hr\hr}\tbr{s}um	&		&	{\hr\hr}la\tbr{s}	&	{\hr\hr}-aa\tbr{s}	&	{\hr\hr}\tbr{s}ita	&	{\hr\hr}-\tbr{s}oo\\
Lola	&	{\hr\hr}\tbr{s}eri	&		&	{\hr\hr}\tbr{s}eman	&	{\hr\hr}la\tbr{s}i	&		&	{\hr\hr}\tbr{h}ita	&	{\hr\hr}-\tbr{s}ewa\\
Lorang	&	{\hr\hr}\tbr{s}era	&	{\hr\hr}ubu	&	{\hr\hr}yemal	&	{\hr\hr}lay\su{Δ}	&	{\hr\hr}-yay	&	{\hr\hr}ʔita	&	\hr\cg{(m-ɸúlar)}\\
Dobel	&	{\hr\hr}yera	&		&	{\hr\hr}yamal	&	{\hr\hr}lay	&		&	{\hr\hr}ʔita	&	\hr\cg{(-yara)}\\
Koba	&	{\hr\hr}\tbr{s}írey	&		&		&	{\hr\hr}lay\su{Δ}	&		&	{\hr\hr}ʔita	&	\\
Barakai	&	{\hr\hr}\tbr{s}er	&	{\hr\hr}\tbr{s}om	&	{\hr\hr}\tbr{s}emal	&	{\hr\hr}laʔ	&		&	{\hr\hr}\tbr{s}e	&	{\hr\hr}-ʤe\\
Karey	&	{\hr\hr}\tbr{s}er	&		&		&	{\hr\hr}lado	&		&	{\hr\hr}\tbr{s}edo	&	\hr\cg{(-wólar)}\\
Batuley	&	{\hr\hr}\tbr{s}er	&	{\hr\hr}\tbr{s}um	&		&	{\hr\hr}lae\tbr{s}\su{Δ}	&	{\hr\hr}-e\tbr{s}	&	{\hr\hr}\tbr{s}it	&	\hr\cg{(it)}\\
Mariri	&	{\hr\hr}\tbr{s}er(i)	&	{\hr\hr}\tbr{s}um	&		&	{\hr\hr}laʤ	&		&	{\hr\hr}\tbr{s}in	&	\\
Manombai	&	{\hr\hr}\tbb{t}era	&	{\hr\hr}\tbr{s}ubu	&	{\hr\hr}\tbb{t}emal	&	{\hr\hr}la\tbr{s}i	&		&	{\hr\hr}\tbr{s}ita	&	\hr\cg{(-erei)}\\
NE. Tarang.	&	{\hr\hr}\tbr{s}ira	&		&		&	{\hr\hr}la\tbb{t}	&		&	{\hr\hr}\tbr{s}êta	&	\hr\cg{(-ʤala)}\\
SE. Tarang.	&	{\hr\hr}\tbr{s}ɛr	&		&		&	{\hr\hr}la\tbb{t}	&		&	{\hr\hr}\tbr{s}ɛr	&	\hr\cg{(-ŋnan)}\\
NW. Tarang.  	&	{\hr\hr}\tbr{s}êra	&		&	{\hr\hr}\tbb{t}emal	&	{\hr\hr}la\tbb{t}	&		&	{\hr\hr}\tbr{s}ita	&	\hr\cg{(-ʤɔla)}\\
SW. Tarang.  	&	{\hr\hr}\tbr{s}êra	&	{\hr\hr}\tbr{s}um	&		&	{\hr\hr}la\tbb{t}	&	{\hr\hr}-e\tbb{t}	&	{\hr\hr}\tbr{s}êta	&	{\hr\hr}-ʤɔw\\
		\lspbottomrule
	\end{tabular}
		\begin{tablenotes}\tnsize
			\item[†]	{Possible cognates of **tera ‘nine’ occur in Flores; e.g.
								Ngadha \it{təra əsa} ‘nine’.}
			\item[‡]	{The full range of meanings for Proto-Aru **-sati is ‘plant
								with dibble, throw, poke/stab’.}
			\item[Δ]	{Lorang \it{lay},
								Koba \it{lay}, and Batuley \it{laes} are \tsc{3pl} forms.
								Other forms are inanimate.}
		\end{tablenotes}
	\end{threeparttable}
\end{adjustbox}
\end{table}

\citet[131-133]{Collins1982c} identifies a \tsc{South Aru} node,
which excludes Kompane and \tsc{Kola}.
There are two exclusively shared innovations which define \tsc{South Aru}.
The first is merger of Proto-Aru **R/**r > \it{r} (see \trf{06-04} and \trf{06-05}).
The second is *p > **ɸ > **β (see \trf{06-11}),
Additional support for \tsc{South Aru} comes from fortition of **w.
This fortition affects most \tsc{South Aru} words,
though there are sporadic exceptions.
Fortition of **w also occurs in Kompane,
and is thus not exclusive to \tsc{South Aru}.
The result of this fortition is usually **w > \it{gw}.
\tsc{West Tarangan} and Manombai have subsequent de-labialisation
**gw > \it{g}, which is also attested in most Karey words
and some Lorang words.
Dobel has **w > **gw > \it{kw}
and Koba has **w > **gw > \it{ʍ},
which was almost certainly from intermediate **kw.
Thus, Dobel and Koba share *w > **gw > **\it{kw}.
Examples of fortition of **w are given in \trf{06-15}.
%with \tsc{South Aru} exceptions indicated.

\begin{table}
\caption[Proto-Aru **w]{Proto-Aru **w\su{†}}\label{tab:06-15}
\begin{adjustbox}{width=\textwidth}
	\begin{threeparttable}\st{2pt}
	\begin{tabular}{ll@{ }l@{ }l@{ }l@{ }l@{ }l@{ }l}\lsptoprule
local gl.	&	{\hr}water	&	{\hr}root	&	{\hr}ySib	&	dry in sun	&	{\hr}rattan	&	wake up	&	sleepy\\
PMP	&	\hr*\tbr{w}ahiyəR	&	\hr*\tbr{w}akaR	&	\hr*huaji	&	\hp{*-}*\tbr{w}aRi	&		&		&	\\
P.-Aru	&	**\tbr{w}ayaR	&	**\tbr{w}akaR	&	**\tbr{w}ali	&	**-\tbr{w}aRay	&	**\tbr{w}aŋguR	&	**-\tbr{w}aŋaR	&	**-\tbr{w}oŋga\\ \midrule
Ujir	&	{\hr\hr}\tbr{w}ay	&	{\hr\hr}\tbr{w}aki	&	{\hr\hr}\tbr{w}ɪl	&	{\hr\hr}-\tbr{w}oy	&	{\hr\hr}\tbr{w}aŋi	&	{\hr\hr}-\tbr{w}aŋay	&	\\
Kompane	&	{\hr\hr}\tbr{gw}ah	&	{\hr\hr}\tbr{w}aki	&	{\hr\hr}\tbr{gw}áel	&	{\hr\hr}-\tbr{gw}ahi	&	{\hr\hr}\tbr{gw}aŋu	&	{\hr\hr}-\tbr{gw}aŋah	&	{\hr\hr}-\tbr{g}oŋa\\
East Kola	&	{\hr\hr}\tbr{w}eeh	&	{\hr\hr}\tbr{w}akih	&	{\hr\hr}\tbr{w}el	&	{\hr\hr}-\tbr{w}eh	&	{\hr\hr}\tbr{w}aguh	&	{\hr\hr}-\tbr{w}aŋah	&	{\hr\hr}-\tbr{w}oga\\
Lola	&	{\hr\hr}\tbr{w}aʔ	&	{\hr\hr}\tbr{w}aʔi	&	{\hr\hr}\tbr{w}en	&	{\hr\hr}-\tbr{w}ahi	&	{\hr\hr}\tbr{w}aŋ	&	{\hr\hr}-\tbr{w}aŋ	&	\\
Lorang	&	{\hr\hr}\tbr{gw}ar	&	{\hr\hr}\tbr{gw}akir	&	{\hr\hr}\tbr{g}el	&	{\hr\hr}-\tbr{g}era	&		&	{\hr\hr}-\tbr{gw}aŋar	&	{\hr\hr}-\tbr{g}oŋa\\
Dobel	&	{\hr\hr}\tbr{kw}ar	&	{\hr\hr}\tbr{kw}aʔir	&	{\hr\hr}\tbr{kw}ali-	&	{\hr\hr}-\tbr{kw}ari	&	{\hr\hr}\tbr{kw}aŋur	&	{\hr\hr}-\tbr{kw}aŋar	&	\\
Koba	&	{\hr\hr}\tbr{ʍ}ar	&	{\hr\hr}\tbr{ʍ}akir	&	{\hr\hr}\tbr{ʍ}éla	&	{\hr\hr}-\tbr{ʍ}era	&	{\hr\hr}\tbr{ʍ}aŋar	&	{\hr\hr}-\tbr{ʍ}áŋar	&	\\
Barakai	&	{\hr\hr}\tbr{gw}aor	&	{\hr\hr}\tbr{gw}áer	&	{\hr\hr}\tbr{gw}el-	&	\hp{*-}\cg{(da-bálar)}	&	{\hr\hr}\tbr{gw}agur\su{Δ}	&		&	{\hr\hr}-\tbr{gw}in\\
Karey	&	{\hr\hr}\tbr{g}aor	&	{\hr\hr}\tbr{g}á\tbr{g}er	&	{\hr\hr}\tbr{g}óliŋ	&	{\hr\hr}-\tbr{ɸ}er	{\cgr}&	{\hr\hr}\tbr{gw}aŋur	&		&	\\
Batuley	&	{\hr\hr}\tbr{gw}ayor	&	{\hr\hr}\tbr{gw}ágwer	&	{\hr\hr}\tbr{gw}al	&	{\hr\hr}\tbr{w}ayer	{\cgr}&	{\hr\hr}\tbr{gw}aŋor	&	{\hr\hr}-\tbr{w}aŋar	{\cgr}&	\\
Mariri	&	{\hr\hr}\tbr{gw}ayor	&	{\hr\hr}\tbr{w}ager	{\cgr}&	{\hr\hr}\tbr{gw}el-	&	{\hr\hr}-\tbr{w}er	&	{\hr\hr}\tbr{gw}aŋor	&		&	\\
Manombai	&	{\hr\hr}\tbr{g}ayar	&	{\hr\hr}\tbr{g}ákar-i	&	{\hr\hr}\tbr{g}ali	&	{\hr\hr}-\tbr{g}aray	&		&	{\hr\hr}-\tbr{g}aŋar	&	{\hr\hr}-\tbr{g}oŋa\\
NE. Tar.	&	{\hr\hr}\tbr{gw}ar	&	{\hr\hr}\tbr{gw}akir	&	{\hr\hr}\tbr{gw}ɛl-	&	{\hr\hr}-\tbr{gw}arɛ	&	{\hr\hr}\tbr{gw}aŋar	&		&	{\hr\hr}-\tbr{g}ɔm\\
SE. Tar.	&	{\hr\hr}\tbr{gw}ar	&	{\hr\hr}\tbr{gw}agwir	&	{\hr\hr}\tbr{gw}ɛl-	&	{\hr\hr}-\tbr{gw}aray	&	{\hr\hr}\tbr{gw}aŋar	&		&	{\hr\hr}-\tbr{g}ɔm{\tl}gɔm\\
NW. Tar.	&	{\hr\hr}\tbr{g}ar	&	{\hr\hr}\tbr{g}akir	&	{\hr\hr}ʤɛl-\su{‡}	&	{\hr\hr}-\tbr{g}aray	&	{\hr\hr}\tbr{g}aŋar	&		&	{\hr\hr}-\tbr{g}ɔm\\
SW. Tar.	&	{\hr\hr}\tbr{g}ar	&	{\hr\hr}\tbr{g}akar	&	{\hr\hr}ʤɛl-	&	{\hr\hr}-\tbr{g}aray	&	{\hr\hr}\tbr{g}aŋar	&	{\hr\hr}-\tbr{g}aŋar	&	{\hr\hr}-\tbr{g}ɔm\\
		\lspbottomrule
	\end{tabular}
		\begin{tablenotes}\tnsize
			\item[†]	{Shading indicates \tsc{South Aru} words which do not have fortition of *w.}
			\item[‡]	{\tsc{West Tarangan} \it{ʤɛl-} ‘younger sibling’ has palatalisation
								of **g before a front vowel.}
			\item[Δ]	{Barakai \it{gwagur} is from Mesiang village.}
		\end{tablenotes}
	\end{threeparttable}
\end{adjustbox}
\end{table}

\citet{Nivens2017} includes Lola within \tsc{South Aru},
but Lola does not have fortition of *w
and does not merge **R/**r > \it{r}.
Instead, Lola retains **r = \it{r}
and has **R > \it{h}.
The change **R > \it{h} also occurs in Kompane
and \tsc{Kola}, and on this basis,
we could place these three languages in a single node.
However, West Kola has *s > \it{h} (see \trf{06-13}), while East
Kola, Kompane, and Lola have *s > (**h) > {\0}.
Thus, if we posited **R > \it{h} at a single node encompassing
these languages, it would incorrectly predict a merger of *j/*R/*s > **R/**h >
{\0} in these languages.
Instead, it seems that the change **R > \it{h} in Lola is independent
of that in \tsc{Kola} and Kompane.
\tsc{Kola} and Kompane are spoken adjacent to one another
and we can posit that **R > \it{h} spread among these languages
by diffusion after East Kola
and Kompane had undergone (*s >) **h > {\0}.
Because no other significant sound changes are shared between
these languages, we place \tsc{Kola}, Kompane,
and Lola as isolates within \tsc{Nuclear Aru}.

Innovation-defined nodes within \tsc{South Aru} are not easy to identify.
Lorang, Dobel, and Koba all share another round of word-initial
glide epenthesis (after the fortition of glides inserted at an
earlier stage), and could be placed together on the basis of this change.
However, this change is also found in Lola which we exclude from \tsc{South Aru}.
Furthermore, glide epenthesis is a recurrent sound change in
the languages of Aru, which has probably been ongoing ever since
the formation of this subgroup (see \srf{06-01-03}).

We identify only three further innovation-defined nodes within \tsc{South Aru}.
Firstly, Dobel and Koba clearly group together on the basis of
*y/*z > **y > **ʤ > \it{s} (see examples in \trf{06-06},
and \trf{06-08}), as well as devoicing of *w > **gw.
Dobel has **gw > \it{kw},
and Koba has **w > **gw > \it{ʍ} (voiceless labio-velar approximant).
This group also matches the lexicostatistical classification
of \citet{Hughes1987} in which Koba was lexically closest to Dobel.

Secondly, Barakai and Karey can be grouped together on the basis
of *b > **p > \it{w} in most words (see \trf{06-11}).
This change is different from most other \tsc{Aru} languages
which have *b > \it{ɸ.} Note,
however, that Karey has *p > \it{ɸ} /{\gap}o.
Again, this classification matches the lexical classification of \citet{Hughes1987}.

Thirdly, we place the \tsc{Tarangan} languages together on the
basis of initial PMP *pa > Proto-Aru **ɸ > Proto-South Aru **βa
> \it{ɔ} in most words.
Two examples of *pa > \it{ɔ} in \tsc{Tarangan} have been given
in \trf{06-11} (*paniki > East Tarangan \it{ɔnít} ‘flying fox’
and *pana > \tsc{West Tarangan} \it{-ɔn} ‘shoot’).
Two additional examples are *panaw ‘whitish skin disorder’ >
Southwest Tarangan \it{ɔnaw} (East Kola \it{ɸʷanaw},
Dobel \it{wanu}), and *paŋdan > Southwest Tarangan \it{ɔdan} ‘pandanus’.
Supporting evidence for \tsc{Tarangan} may come from
the innovation of seven-vowel systems in which mid-high ‹ê›
/e/ and ‹ô› /o/ contrast with mid-low /ɛ/ and /ɔ/.

No further innovation-defined nodes are well supported within
\tsc{South Aru} on the basis of the phonological history.
\citet{Hughes1987} treated Mariri as a dialect of Batuley,
but we find little support for this from the historical phonology.
They do share medial *-k- > \it{-g-},
but this change is also found in Karey.
Other changes that are shared by Batuley
and Mariri, such as *p > {\0},
are common among \tsc{Aru} languages.

One unusual change which deserves mention occurs in Mariri,
which has *n > \it{d} word initially and medially.
Examples include: *anak > \it{adeʤ} ‘child’,
*naŋi > \it{-den} ‘swim’,
*nipən > \it{deen} ‘tooth’,
*bunuq > \it{-ɸʊdʊn} ‘kill’,
*kunij > \it{kódel} ‘yellow’,
*niuR > \it{dor} ‘coconut’,
*panua > \it{ɸadú} ‘village’,
and *haŋin > **anin > \it{aden} ‘wind’.

